% =========================================================
% T3 — RELATIONAL GEOMETRY (Hyper‑Condensed)
% =========================================================

\section{T3 Relational Geometry}

\paragraph{Overview.}
Relational Geometry defines the lawful structure governing cross-system
interaction. It provides the primitives, fields, invariants, drift geometry,
legality constraints, and dual-system coupling rules that allow human and
machine systems to interact without violating curvature, identity, or legality
conditions.

% ---------------------------------------------------------
% T3.1 Relational Primitives
% ---------------------------------------------------------
\subsection{T3.1 Relational Primitives}
Relational primitives define the atomic units of relational interaction:
relational substrates, relational links, relational signals, relational
invariants, relational drift forms, and relational legality primitives. These
primitives form the base geometry for all cross-system coupling.

% ---------------------------------------------------------
% T3.2 Relational Fields
% ---------------------------------------------------------
\subsection{T3.2 Relational Fields}
Relational fields define the lawful regions where relational primitives
interact. Field envelopes, field channels, field stability modes, and field
legality surfaces govern how signals, invariants, drift, and legality propagate
across relational space.

% ---------------------------------------------------------
% T3.3 Relational Invariants
% ---------------------------------------------------------
\subsection{T3.3 Relational Invariants}
Relational invariants define the stability conditions preserved across
relational transitions. Meaning invariants, identity invariants, geometry
invariants, and legality invariants ensure that relational interaction remains
lawful under drift, load, and deformation.

% ---------------------------------------------------------
% T3.4 Relational Drift Geometry
% ---------------------------------------------------------
\subsection{T3.4 Relational Drift Geometry}
Relational drift geometry defines how drift propagates across relational
surfaces. Drift channels, drift envelopes, drift loops, and cross-field drift
propagation describe lawful deformation under relational load and cross-system
pressure.

% ---------------------------------------------------------
% T3.5 Relational Legality Geometry
% ---------------------------------------------------------
\subsection{T3.5 Relational Legality Geometry}
Relational legality geometry defines legality constraints for relational
coupling. Legality masks, legality channels, legality propagation rules, and
legality-preserving operators ensure that relational transitions remain lawful
across human, machine, and federated systems.

% ---------------------------------------------------------
% T3.6 Human-Machine Relational Dynamics
% ---------------------------------------------------------
\subsection{T3.6 Human-Machine Relational Dynamics}
Human-machine relational dynamics define dual-lens relational thermodynamics.
Human curvature (agency, teleology, identity) and machine curvature (deployment,
execution, runtime) interact through relational fields, producing shared drift
channels, collapse trajectories, and recovery envelopes. Dual-system stability
requires lawful relational coupling across all manifold layers.
